% ============================================================
% CAPÍTULO 7: VITAMINAS
% ============================================================
% CONTENIDO BASADO EN: Unidad de Aprendizaje II, Tema "Vitaminas"
% PROPÓSITO: Describir la composición química, estructura, clasificación y función de las vitaminas, y su importancia biotecnológica.
% ============================================================

\section{Composición química y estructura de las vitaminas}

% --- CONSEJO: Define qué son las vitaminas: compuestos orgánicos esenciales que el cuerpo no puede sintetizar en cantidades suficientes y deben obtenerse de la dieta. ---
\resumebox{definicion}{Vitaminas}{Son compuestos orgánicos heterogéneos, esenciales para el metabolismo celular. Actúan principalmente como coenzimas o precursores de coenzimas.}

% --- CONSEJO: Menciona que las vitaminas no son una fuente de energía, pero son cruciales para el correcto funcionamiento del metabolismo. ---

\section{Clasificación de las vitaminas}

% --- CONSEJO: La clasificación más común es en función de su solubilidad. ---
\begin{table}[H]
    \centering
    \caption{Clasificación de las vitaminas}
    \label{tab:clasificacion_vitaminas}
    \begin{tabular}{|p{0.25\linewidth}|p{0.35\linewidth}|p{0.3\linewidth}|}
        \hline
        \theader{Clasificación} & \theader{Vitaminas} & \theader{Características} \\
        \hline
        \multirow{4}{0.25\linewidth}{\centering \textbf{Liposolubles}} & Vitamina A (Retinol) & Se almacenan en el hígado y tejido adiposo. \\
        & Vitamina D (Colecalciferol) & Se absorben con las grasas de la dieta. \\
        & Vitamina E (Tocoferoles) & Su exceso puede ser tóxico (hipervitaminosis). \\
        & Vitamina K (Filoquinonas, Menaquinonas) & \\
        \hline
        \multirow{9}{0.25\linewidth}{\centering \textbf{Hidrosolubles}} & Vitamina C (Ácido ascórbico) & No se almacenan en el cuerpo (se excretan por orina). \\
        & Vitamina B\textsubscript{1} (Tiamina) & Se absorben fácilmente. \\
        & Vitamina B\textsubscript{2} (Riboflavina) & Su exceso generalmente se excreta (riesgo bajo de toxicidad). \\
        & Vitamina B\textsubscript{3} (Niacina) & \\
        & Vitamina B\textsubscript{5} (Ácido pantoténico) & \\
        & Vitamina B\textsubscript{6} (Piridoxina) & \\
        & Vitamina B\textsubscript{7} (Biotina) & \\
        & Vitamina B\textsubscript{9} (Ácido fólico) & \\
        & Vitamina B\textsubscript{12} (Cobalamina) & \\
        \hline
    \end{tabular}
\end{table}

\section{Función de las vitaminas a nivel celular}

% --- CONSEJO: Explica el papel principal de las vitaminas como coenzimas. ---
\begin{itemize}
    \item \textbf{Coenzimas:} Muchas vitaminas hidrosolubles son precursores de coenzimas, moléculas orgánicas que se unen a las enzimas y son necesarias para su actividad catalítica.
    \item \textbf{Antioxidantes:} Vitaminas como la E y la C protegen a las células del daño oxidativo.
    \item \textbf{Regulación hormonal y de la expresión génica:} La vitamina A y la D actúan como hormonas o precursores de hormonas, regulando la expresión de genes.
\end{itemize}

% --- CONSEJO: Puedes incluir una tabla con ejemplos de vitaminas, sus coenzimas derivadas y las reacciones que catalizan. ---

\section{Importancia biotecnológica de las vitaminas}

% --- CONSEJO: Esta sección es clave para la asignatura. Explica cómo la biotecnología se relaciona con la producción y aplicación de vitaminas. ---
\begin{itemize}
    \item \textbf{Producción industrial:} La biotecnología permite la producción de vitaminas mediante fermentación microbiana, un proceso más sostenible y eficiente que la síntesis química. Ejemplos: producción de vitamina B\textsubscript{12} y B\textsubscript{2} (riboflavina) con bacterias y hongos.
    \item \textbf{Industria alimentaria:} Uso de vitaminas como aditivos para enriquecer alimentos (biofortificación).
    \item \textbf{Industria farmacéutica:} Producción de suplementos vitamínicos y medicamentos para tratar deficiencias.
    \item \textbf{Ingeniería metabólica:} Modificación genética de microorganismos para mejorar la producción de vitaminas.
\end{itemize}